\chapter{Lemma 11.5 and 11.6 — AC composition}
\label{chap:11_5}

Reference: [BMSZ] Lemma~11.5 (two-step AC formula) and Lemma~11.6 (first-entry
of $\mathcal{L}_\tau$); internal blueprints
\texttt{lemma\_11\_5\_PBP\_identity.md}, \texttt{lemma\_11\_5\_COMPLETE.md}.

\section{Setup for Lemma 11.5}

Let $\star \in \{B, D\}$ and assume $\mathbf{r}_2(\check{\mathcal{O}}) > \mathbf{r}_3(\check{\mathcal{O}})$
(the \emph{primitive case}). Set:
\begin{itemize}
  \item $\tau \in \PBP_\star(\check{\mathcal{O}})$ (outer PBP).
  \item $\tau' = \nabla\tau \in \PBP_{\star'}(\check{\mathcal{O}}')$ (single descent; $\star' = C$ for $\star = D$, $\star' = \tilde{C}$ for $\star \in \{B^\pm\}$).
  \item $\tau'' = \nabla^2\tau = \nabla(\tau') \in \PBP_\star(\check{\mathcal{O}}'')$ (double descent).
  \item $\tau_t \in \PBP_D(\check{\mathcal{O}}_t)$ (tail PBP).
  \item $n_0 := (\mathbf{c}_1(\mathcal{O}') - \mathbf{c}_2(\mathcal{O}')) / 2$ (the inner augmentation parameter).
  \item $\mathcal{L}_\tau, \mathcal{L}_{\tau''}$ denote the AC results.
\end{itemize}

By the shape relations $\mathbf{c}_i(\mathcal{O}') = \mathbf{c}_{i+1}(\mathcal{O})$, we have
$n_0 = (\mathbf{c}_2(\mathcal{O}) - \mathbf{c}_3(\mathcal{O}))/2$.

\section{Lemma 11.5 — PBP identity}

\begin{lemma}[Lemma 11.5 D (primitive case, ILS-level, unconditional)]
  \label{lem:11_5_D_ILS}
  \lean{ILS.lemma_11_5_D_unconditional, prop_11_5_D_unconditional, ILS.lemma_11_5_D, prop_11_5_D}
  \leanok
  \uses{prop:11_4, def:AC_step, def:charTwistCM, def:augment, def:twistBD, thm:AC_step_fcSign, thm:fcSign_diff}
  Given an ILS $E$ (intended as $\mathcal{L}_{\tau''}$) with the firstColSign
  invariant
  \[ \mathrm{fcSign}(E) = (\Sign(E).1 - n_{\mathrm{prev}},\ \Sign(E).2 - n_{\mathrm{prev}}) \]
  and target tail signature $(p_t, q_t)$ satisfying
  \[ p_\tau - 2 n_{\mathrm{inner}} + \Sign(E).2 = p_t,\quad q_\tau - 2 n_{\mathrm{inner}} + \Sign(E).1 = q_t, \]
  define $n_0 := n_{\mathrm{inner}} - \Sign(E).1 - \Sign(E).2 + n_{\mathrm{prev}}$ and
  $E_{\mathrm{inner}} := T^{\gamma_1}\big(E \oplus (n_0, n_0)\big)$. Then:
  \begin{align*}
    p_\tau - \Sign(E_{\mathrm{inner}}).1 - \mathrm{fcSign}(E_{\mathrm{inner}}).2 &= p_t, \\
    q_\tau - \Sign(E_{\mathrm{inner}}).2 - \mathrm{fcSign}(E_{\mathrm{inner}}).1 &= q_t.
  \end{align*}
  These are the ``outer augment'' parameters consumed by the outer C$\to$D or M$\to$B lift.
\end{lemma}

\begin{proof}
  \uses{lem:augment_sign, lem:twist_sign, lem:twist_fcSign, thm:fcSign_diff, prop:11_4, lem:sig_sum_D, lem:residual_D}
  \leanok
  \textbf{Step 1 — unfold sign and fcSign of $E_{\mathrm{inner}}$.}
  $T^\gamma$ preserves both sign and fcSign (Lem~\ref{lem:twist_sign}, \ref{lem:twist_fcSign}), so we may work with $F := E \oplus (n_0, n_0) = (n_0, n_0) :: E$.
  By Lem~\ref{lem:augment_sign} (for $n_0 \geq 0$):
  \begin{align*}
    \Sign(F) &= (n_0 + \Sign(E).1 + \mathrm{fcSign}(E).2,\ n_0 + \Sign(E).2 + \mathrm{fcSign}(E).1). \\
    \mathrm{fcSign}(F) &= (n_0 + \mathrm{fcSign}(E).2,\ n_0 + \mathrm{fcSign}(E).1)
  \end{align*}
  using the cons formula for fcSign (the swap is due to parity index shift, Eq.~(9.12)).

  Substituting $\mathrm{fcSign}(E) = (\Sign(E).1 - n_{\mathrm{prev}},\ \Sign(E).2 - n_{\mathrm{prev}})$:
  \begin{align*}
    \Sign(F).1 &= n_0 + \Sign(E).1 + \Sign(E).2 - n_{\mathrm{prev}} = n_{\mathrm{inner}} \\
    \mathrm{fcSign}(F).2 &= n_0 + \Sign(E).1 - n_{\mathrm{prev}} = n_{\mathrm{inner}} - \Sign(E).2.
  \end{align*}
  where we used the definition $n_0 = n_{\mathrm{inner}} - \Sign(E).1 - \Sign(E).2 + n_{\mathrm{prev}}$.

  \textbf{Step 2 — substitute into the outer augment formula.}
  $p_\tau - \Sign(F).1 - \mathrm{fcSign}(F).2 = p_\tau - n_{\mathrm{inner}} - (n_{\mathrm{inner}} - \Sign(E).2) = p_\tau - 2 n_{\mathrm{inner}} + \Sign(E).2 = p_t$ (by hypothesis).

  The $q$-equation is symmetric.
\end{proof}

\section{Lemma 11.5 atomic PBP discharge}

The ILS-level Lemma~\ref{lem:11_5_D_ILS} takes $E$ and the hypotheses as parameters.
At the PBP level, three of the hypotheses become atomic cell-count facts that
are \emph{structurally} true:

\begin{lemma}[Fact 1 — signature sum D]
  \label{lem:fact1_sig_sum}
  \lean{PBP.signature_sum_D, PaintedYoungDiagram.card_eq_sum_countSym}
  \leanok
  \uses{def:signature_pbp, def:cardCells}
  $p_\tau + q_\tau = 2|\tau|$.
\end{lemma}

\begin{lemma}[Fact 2 — descent cell-count definitional]
  \label{lem:fact2_descent}
  \leanok

  \uses{def:cardCells, def:doubleDescent_D}
  $|\tau| - \mathbf{c}_1 = $ single-descent cell count (by definition of cells - col 0).
\end{lemma}

\begin{lemma}[Fact 3 — tail residual]
  \label{lem:fact3_residual}
  \lean{PBP.residual_identity_D}
  \leanok
  \uses{def:tailSignature}
  $p_{\tau_t} + q_{\tau_t} = 2 \mathbf{c}_1 - 2 \mathbf{c}_2$.
\end{lemma}

\begin{theorem}[Lemma 11.5 atomic discharge at PBP level]
  \label{thm:11_5_atomic_pbp}
  \lean{prop_11_5_D_atomic, prop_11_5_D_atomic_pbp_discharged}
  \leanok
  \uses{lem:11_5_D_ILS, lem:fact1_sig_sum, lem:fact2_descent, lem:fact3_residual}
  Given the canonical choices of $|\tau|, \mathbf{c}_1, \mathbf{c}_2,$ tail signatures,
  the three atomic cell-count hypotheses of \texttt{prop\_11\_5\_D\_atomic} are
  discharged automatically; the remaining inputs are:
  \begin{itemize}
    \item $(E, n_{\mathrm{prev}}, p_{\mathrm{prev}}, q_{\mathrm{prev}})$: an inner ILS + its Sign and prev-sig.
    \item $h_{\mathrm{fc}}, h_{\mathrm{sign}}$: firstColSign/Sign relations on $E$ (auto by AC.fold).
    \item $h_{\mathrm{prop11.4}}$: signature decomposition (Prop~\ref{prop:11_4}).
    \item $h_{n_0}$: $n_0 \geq 0$.
  \end{itemize}
\end{theorem}

\section{Lemma 11.6 — first entry of $\mathcal{L}_\tau$}

\begin{lemma}[Lemma 11.6 D (AC-level, unconditional)]
  \label{lem:11_6_D}
  \lean{AC.lemma_11_6_D_unconditional}
  \leanok
  \uses{lem:11_5_D_ILS, thm:AC_step_sign_D, thm:AC_step_fcSign, thm:postTwist_first_entry}
  Let $\gamma = D$ (no C/M pre-twist). For \texttt{source : ACResult} with uniform
  $\Sign = \mathrm{src\_sig}$ and uniform $\mathrm{fcSign} = \mathrm{src\_fcSig}$ satisfying the
  standard-case condition ($p - \mathrm{src\_sig}.1 - \mathrm{src\_fcSig}.2 \geq 0$, etc.):
  every branch $r \in \mathrm{AC.step}(\mathrm{source}, D, p, q, \eps_\tau, \eps_\wp)$ has
  \[ r.2.\mathrm{head?} = \mathrm{some}\big(p - \mathrm{src\_sig}.1 - \mathrm{src\_fcSig}.2,\ (-1)^{\eps_\tau}(q - \mathrm{src\_sig}.2 - \mathrm{src\_fcSig}.1)\big). \]
\end{lemma}

\begin{proof}
  \leanok
  By Thm~\ref{thm:thetaLift_first_entry}, the first entry of $\hat\vartheta_{CD}(\mathrm{source})$ (standard case) is
  $(p - \mathrm{src\_sig}.1 - \mathrm{src\_fcSig}.2,\ q - \mathrm{src\_sig}.2 - \mathrm{src\_fcSig}.1)$.
  The post-twist $\otimes(1, -1)$ (when $\eps_\tau = 1$) maps $(a, b) \mapsto (a, -b)$ by Thm~\ref{thm:postTwist_first_entry}.
\end{proof}

\begin{lemma}[Lemma 11.6 $B^+ / B^-$]
  \label{lem:11_6_B}
  \lean{AC.lemma_11_6_Bplus_unconditional, AC.lemma_11_6_Bminus_unconditional}
  \leanok
  \uses{lem:11_5_D_ILS}
  Parallel statements for $\gamma \in \{B^+, B^-\}$, using $\hat\vartheta_{MB}$.
\end{lemma}

\begin{corollary}[PBP-level Lemma 11.6]
  \label{cor:11_6_pbp}
  \lean{AC.lemma_11_6_D_unconditional}
  \leanok
  \uses{lem:11_6_D, prop:11_4}
  At the PBP level, when $\mathrm{src\_sig} = \Sign(\tau'') = (p_{\tau''}, q_{\tau''})$
  and $\mathrm{src\_fcSig} = \Sign(\tau'') - \Sign(\tau^{(4)})$, the first entry of
  $\mathcal{L}_\tau$ equals
  \[ \big(p_{\tau_t},\ (-1)^{\eps_\tau} q_{\tau_t}\big) \]
  by combining Lem~\ref{lem:11_6_D} with Prop~\ref{prop:11_4} (which gives $p_\tau - \mathbf{c}_2 - p_{\tau''} = p_{\tau_t}$).
\end{corollary}
